L'étiquette d'une bouteille de vinaigre commercial indique $6^{\circ}$. La
concentration molaire en acide éthanoïque dans ce vinaigre est $C_{0}$.

On se propose de doser par pH-métrie ce vinaigre afin de déterminer son degré
d'acidité. Pour cela, on prépare une solution aqueuse $\mathrm{(S_{1})}$ par
dilution 10 fois du vinaigre commercial et on prélève un volume
${V_{A}=\qty{25}{\mL}}$ de la solution diluée $\mathrm{(S_{1})}$ de concentration
molaire $C_{A}$ $\left(C_{A}=\frac{C_{0}}{10}\right)$ que l'on dose avec une
solution aqueuse $\mathrm{(S_{2})}$ d'hydroxyde de sodium
$\ce{Na^{+}_{(aq)} + HO^{-}_{(aq)}}$ de concentration molaire
$C_{B}=\qty{2.5e-1}{\mol\per\L}$.

À l'équivalence, le volume de la solution $\mathrm{(S_{2})}$ ajouté est
$V_{B,E}=\qty{10}{\mL}$.

\begin{enumerate}
    \item Écrire l'équation de la réaction qui a eu lieu lors du dosage,
          supposée totale.
    \item Calculer la valeur de $C_{A}$. En déduire la valeur de $C_{0}$.
    \item Vérifier la valeur du degré d'acidité du vinaigre indiquée sur
          l'étiquette de la bouteille.
\end{enumerate}
